\textbf{R:} Etiquetamos los puntos: $\mathbf{x}^{(1)} = (0,0)$, $\mathbf{x}^{(2)} = (0.2, 0.2)$, $\mathbf{x}^{(3)} = (0.7, 0.7)$, $\mathbf{x}^{(4)} = (0.5, 0.5)$, $\mathbf{x}^{(5)} = (1, 1)$.\\

\textbf{Obs:} Todos los puntos están sobre la recta $x_1 = x_2$. Para $\mathbf{x} = (a, a)$ y un centroide $\boldsymbol{\mu} = (c, c)$,
\[
    \| \mathbf{x} - \boldsymbol{\mu} \|^2 = (a-c)^2 + (a-c)^2 = 2 (a-c)^2.
\]
Usaremos la distancia al cuadrado (asignar por la menor $\|\cdot\|^2$ es lo mismo que asignar por la menor distancia).\\

El algoritmo alterna dos pasos: \textbf{(A) Asignación} de cada punto al centroide más cercano y \textbf{(B) Actualización} de cada centroide como la media de los puntos asignados.

\subsection*{Iteración 1}
Centroides: $\boldsymbol{\mu}_1 = (0.1, 0.1)$, $\boldsymbol{\mu}_2 = (0.6, 0.6)$.

\subsubsection*{(A) Asignación}
Distancias al cuadrado $d^2 = 2(a-c)^2$:
\begin{center}
    \begin{tabular}{ccccc}
        \toprule
        Punto & $\| \mathbf{x} - \boldsymbol{\mu}_1 \|^2$ & $\| \mathbf{x} - \boldsymbol{\mu}_2 \|^2$ & Más cercano & Clúster \\
        \midrule
        $(0, 0)$     & $2(0.1)^2 = 0.02$ & $2(0.6)^2 = 0.72$ & $\boldsymbol{\mu}_1$ & $C_1$ \\
        $(0.2, 0.2)$ & $2(0.1)^2 = 0.02$ & $2(0.4)^2 = 0.32$ & $\boldsymbol{\mu}_1$ & $C_1$ \\
        $(0.7, 0.7)$ & $2(0.6)^2 = 0.72$ & $2(0.1)^2 = 0.02$ & $\boldsymbol{\mu}_2$ & $C_2$ \\
        $(0.5, 0.5)$ & $2(0.4)^2 = 0.32$ & $2(0.1)^2 = 0.02$ & $\boldsymbol{\mu}_2$ & $C_2$ \\
        $(1, 1)$     & $2(0.9)^2 = 1.62$ & $2(0.4)^2 = 0.32$ & $\boldsymbol{\mu}_2$ & $C_2$ \\
        \bottomrule
    \end{tabular}
\end{center}
Es decir,
\[
    C_1 = \{(0, 0), (0.2, 0.2)\}, \qquad C_2 = \{(0.7, 0.7), (0.5, 0.5), (1, 1)\}.
\]

\subsubsection*{(B) Actualización}
\begin{align*}
    \boldsymbol{\mu}_1 &= \tfrac{1}{2} [(0, 0) + (0.2, 0.2)] = \big(\tfrac{0.2}{2}, \tfrac{0.2}{2}\big) = (0.1, 0.1), \\
    \boldsymbol{\mu}_2 &= \tfrac{1}{3} [(0.7, 0.7) + (0.5, 0.5) + (1, 1)] = \big(\tfrac{2.2}{3}, \tfrac{2.2}{3}\big) \approx (0.7333, 0.7333).
\end{align*}
$\boldsymbol{\mu}_1$ no se mueve (la media de su clúster coincide con su valor inicial); $\boldsymbol{\mu}_2$ pasa de $0.6$ a $\approx 0.7333$.

\subsection*{Iteración 2}
Centroides: $\boldsymbol{\mu}_1 = (0.1, 0.1)$, $\boldsymbol{\mu}_2 \approx (0.7333, 0.7333)$.

\subsubsection*{(A) Asignación}
\begin{center}
    \begin{tabular}{ccccc}
        \toprule
        Punto & $\| \mathbf{x} - \boldsymbol{\mu}_1 \|^2$ & $\| \mathbf{x} - \boldsymbol{\mu}_2 \|^2$ & Más cercano & Clúster \\
        \midrule
        $(0, 0)$     & $0.0200$ & $2(0.7333)^2 \approx 1.0756$ & $\boldsymbol{\mu}_1$ & $C_1$ \\
        $(0.2, 0.2)$ & $0.0200$ & $2(0.5333)^2 \approx 0.5689$ & $\boldsymbol{\mu}_1$ & $C_1$ \\
        $(0.7, 0.7)$ & $0.7200$ & $2(0.0333)^2 \approx 0.0022$ & $\boldsymbol{\mu}_2$ & $C_2$ \\
        $(0.5, 0.5)$ & $0.3200$ & $2(0.2333)^2 \approx 0.1089$ & $\boldsymbol{\mu}_2$ & $C_2$ \\
        $(1, 1)$     & $1.6200$ & $2(0.2667)^2 \approx 0.1422$ & $\boldsymbol{\mu}_2$ & $C_2$ \\
        \bottomrule
    \end{tabular}
\end{center}
Las asignaciones no cambian respecto a la iteración 1:
\[
    C_1 = \{(0, 0), (0.2, 0.2)\}, \qquad C_2 = \{(0.5, 0.5), (0.7, 0.7), (1, 1)\}.
\]

\subsubsection*{(B) Actualización}
Al recalcular las medias se obtienen los mismos centroides, $\boldsymbol{\mu}_1 = (0.1, 0.1)$ y $\boldsymbol{\mu}_2 \approx (0.7333, 0.7333)$, así que el algoritmo ha convergido (ni las asignaciones ni los centroides cambian).

\subsection*{Resultado final}
\begin{align*}
    C_1 &= \{(0, 0), (0.2, 0.2)\}, & \boldsymbol{\mu}_1 &= (0.1, 0.1), \\
    C_2 &= \{(0.5, 0.5), (0.7, 0.7), (1, 1)\}, & \boldsymbol{\mu}_2 &\approx (0.7333, 0.7333).
\end{align*}
Valor final de la función objetivo (SSE):
\[
    J = \underbrace{(0.02 + 0.02)}_{C_1} + \underbrace{(0.0022 + 0.1089 + 0.1422)}_{C_2} \approx 0.0400 + 0.2533 = 0.2933.
\]
La separación natural de los datos en un grupo ``bajo'' $\{0, 0.2\}$ y uno ``alto'' $\{0.5, 0.7, 1\}$ se obtiene ya desde la primera iteración.
